% SPDX-License-Identifier: Apache-2.0
\section{An Offer That Does Not Wait for Room}
\label{sec:x-put-axi}

A \code{put} waits for room before it raises \code{valid}, so its
\code{valid} waits on its \code{ready}.
Inside TxHDL that is harmless, since every join goes through a channel
whose \code{ready} is a register.
AXI forbids it of a source, though, and a sender that waits on
\code{ready} joined wire to wire to a receiver that waits on
\code{valid} passes nothing (\#1868).
\code{\#[axi(a, b)]} on \code{run} marks sending ports whose
\code{put} offers from the first edge it waits at, without waiting for
room (\#1871).
Its \code{valid} is the state's condition alone, low while the reset
is high, as AXI also requires; the state is still left on
\code{ready}, so what follows the \code{put} happens on the take, as
before.
The word is read at the first edge of the wait and held until it is
taken, since AXI requires the data to hold while it waits; in the
netlist that is two registers per marked port, \code{<port>\_axi\_wait}
and \code{<port>\_axi\_hold}, and no cycle of latency.
A port the attribute does not name is as in the last section.
A marked port sends with \code{put} alone, since a \code{send} offers
only once there is room, and \code{\#[axi]} on a receiving port is
refused, since a receiver's \code{ready} may wait on \code{valid}
under AXI.

\code{Offers} has two marked ports.
On \code{a} it puts its count in a loop of one wait; on \code{b} it
puts a pair in a sequence of three waits.
Every word mixes in \code{seed}, an input the run changes every cycle,
and two consumers take when bits of a small shift register say so.
So a word offered across a stall is the one read at its first edge,
and the second word of each pair on \code{b}, which no input touches,
counts up, as the run asserts.
The netlist is checked against the run under nvc and Verilator, and
its header names the two ports as AXI's.
In \code{//lib/chanloop}, the same counter with its port marked passes
a word a cycle joined wire to wire to a stage whose \code{ready} waits
on \code{valid}, where the unmarked one passes none.

\begin{figure}[htbp]
\centering
\input{axi_put_timing.tex}
\caption{Two marked ports: \code{a\_valid} and \code{b\_valid} up from
the first edge of each wait, each word held while \code{a\_ready} or
\code{b\_ready} is low, though \code{seed} changes every cycle.}
\label{fig:x-put-axi}
\end{figure}

\lstinputlisting[caption={\code{ex\_axi\_put.rs}. Run with \code{bazel run //lib/examples:ex\_axi\_put}.},label={lst:x-put-axi}]{lib/examples/ex_axi_put.rs}

\lstinputlisting[style=ascii,caption={What \code{ex\_axi\_put} printed when the build ran it: \code{seed} and each word as it was taken, then the netlist.},label={lst:x-put-axi-out}]{out_axi_put.txt}
